\section{Requirements and Message Protocol}
\label{sec:protocol}

The link carries printable ASCII, one message per line. The client sends a
request and the server answers with exactly one line. The server never speaks
first, so the PC side never has to poll for unsolicited traffic.

\subsection{One newline-terminated line bounds every message}
A message is terminated by \texttt{\textbackslash n}. A
\texttt{\textbackslash r} immediately before it is accepted and discarded, so a
terminal that sends CRLF interoperates without a separate mode. Four limits
bound what the parser will accept:

\begin{itemize}
  \item a frame is at most 64 bytes including the terminator, which leaves 63
        characters of payload; a longer line is discarded up to and including
        the next \texttt{\textbackslash n} and answered
        \texttt{NACK:TOOLONG};
  \item a frame left incomplete for 2000\,ms is cleared, so a client that dies
        mid-line cannot wedge the parser;
  \item bytes outside printable ASCII, that is below \texttt{0x20} or above
        \texttt{0x7E}, are dropped as they arrive rather than buffered;
  \item commands are uppercase, and anything unrecognised is
        \texttt{NACK:BADCMD}.
\end{itemize}

None of these paths can change the game. A rejected, truncated or malformed
request leaves the board, the turn and the game status exactly as they were.

\subsection{The board travels as nine characters}
The board is nine characters, index 0 to 8, row major, where \texttt{.} is an
empty cell and \texttt{X} and \texttt{O} are the players. X always moves first,
so the client derives whose turn it is by counting: equal counts means X to
move. The client therefore holds no authoritative state of its own.

\subsection{Every request has exactly one documented reply}
Table~\ref{tab:protocol} lists the command set. The reader should notice that
every rejection is named: there is no request whose failure produces silence,
which is what lets the client distinguish a refused request from a lost one.

\begin{table}[htbp]
\caption{Command set and responses}
\label{tab:protocol}
\centering
\footnotesize
\begin{tabular}{p{1.8cm}p{3.1cm}p{2.0cm}}
\hline
\textbf{Request} & \textbf{Response} & \textbf{Meaning} \\
\hline
\texttt{CONNECT} & \texttt{OK} & open a session \\
\texttt{NEW} & \texttt{OK:BOARD:.........} & start a game, X to move \\
\texttt{MOVE:<i>} & \texttt{OK:BOARD:<9>} & legal move \\
\texttt{MOVE:<i>} & \texttt{WIN:X}, \texttt{WIN:O}, \texttt{DRAW} & the move ended it \\
\texttt{MOVE:<i>} & \texttt{NACK:RANGE}, \texttt{NACK:OCCUPIED}, \texttt{NACK:NOTURN} & rejected, nothing changed \\
\texttt{STATE} & \texttt{BOARD:<9>} & current board \\
\texttt{RESIGN} & \texttt{OK:GAMEOVER} & forfeit \\
\texttt{DISCONNECT} & \texttt{OK} & end session \\
\hline
\end{tabular}
\end{table}

\subsection{Four departures from the reference protocol}
Each departure closes a failure the reference protocol in the brief leaves
open, and none of them changes a response the brief specifies.

\textbf{\texttt{NACK:TOOLONG}} exists because the reference defines no reply
for an over-length line, which would otherwise leave the client waiting on a
response that never comes.

\textbf{The optional player field on \texttt{MOVE}}, written
\texttt{MOVE:<i>:<p>} where \texttt{<p>} is \texttt{X} or \texttt{O}, lets a
client state whose move it believes it is sending. The server then answers
\texttt{NACK:NOTURN} rather than silently applying the move to the player whose
turn it actually is. Without such a field a single client cannot be out of
turn, so the brief's own out-of-turn requirement would be untestable.

\textbf{\texttt{DISCONNECT} from IDLE} is accepted so that a client
reconnecting after a crash does not have to know what state it left behind. It
is still refused mid-game.

\textbf{An optional \texttt{*HH} checksum suffix} may be appended to any
request. Its scope and its limits are discussed in
Section~\ref{sec:robust}.
